\section{Introduction}
\label{sec:intro}

Figures carry much of the evidence in a scientific paper, but a figure cannot be understood from its pixels
alone. The text explains why the experiment was done, what the figure shows, and what the authors conclude from
it. Collecting that text for each figure is useful in itself (for reading tools, figure search and figure
summarisation), and it is the first step towards the longer-term aim that motivated this project: a dataset of
\{figure, text that discusses it\} pairs for fine-tuning an open vision--language model.

\paragraph{The task.} Given a directory of research papers, output for every figure all the sentences that talk
about it. These sentences are of two kinds.
\begin{itemize}
  \item An \textbf{explicit reference} names the figure: ``Fig.~2a'', ``Figures 3--5'', ``(4B)''.
  \item An \textbf{implicit reference} discusses the content of the figure without naming it. It may sit next
        to an explicit reference, or several pages away.
\end{itemize}
Figure~\ref{fig:example} shows a real example from our benchmark with all three situations that the rest of
this report distinguishes: explicit, implicit \emph{near} a citation, and implicit \emph{far} from any citation.

\begin{figure}[h]
\centering
\fcolorbox{navy!30}{soft}{\begin{minipage}{0.95\linewidth}\small
\textbf{Figure 2 of the article} (caption, shortened): \emph{Production of interleukin (IL)-17 by peripheral blood
mononuclear cells \ldots}\\[5pt]
\textcolor{steel}{\textbf{Near implicit} (sentence 107, immediately before the citation).}
To evaluate the effects of inflammatory cytokines released by activated PBMC, we tested the effects of several
cytokines and chemokines on IL-17 production.\\[4pt]
\textcolor{navy}{\textbf{Explicit} (sentences 108--109).}
We detected an increase in IL-17 level after stimulation with IL-15 (10 ng/ml), whereas with IL-1$\beta$ \ldots{}
the levels in IL-17 were unchanged (Fig.~2a).\\[4pt]
\textcolor{accent}{\textbf{Far implicit} (sentence 150, in the Discussion, about forty sentences later).}
Our studies demonstrated that IL-15 and MCP-1 as well as TCR ligation significantly increased the production of
IL-17 in RA PBMC.
\end{minipage}}
\caption{Human gold labels for one figure in \figsum{} (article 1064895, Figure 2). The annotators marked seven
sentences for this figure; only two of them name it.}
\label{fig:example}
\end{figure}

\paragraph{Why it is hard.} Explicit references are a pattern-matching problem and can be solved almost
completely. Implicit references need an understanding of what the figure shows and of how the paper's argument
flows. They are also subjective: two trained biologists labelling the same articles agree only at Cohen's
$\kappa = 0.68$ \citep{ramesh2015figsum}. Finally, labelled data is extremely scarce. We found a single public
benchmark, with 18 usable articles.

\paragraph{Contributions.}
\begin{enumerate}
  \item \textbf{A working tool.} A command-line program that parses PDFs or structured text, detects figures and
        captions, finds explicit references with a grammar, finds implicit references, and writes a JSON file
        and an HTML report per paper (Section~\ref{sec:pipeline}).
  \item \textbf{A strong rule baseline.} \bof{}: the sentence before and after each citation, in the same
        paragraph, minus three kinds of off-target sentences. It scores F1 0.752 on our training papers, against
        0.659 for the best published method on the same papers (Section~\ref{sec:rules}).
  \item \textbf{Learned models that beat the rule on unseen papers.} A cross-encoder and a whole-paper span
        tagger, combined with an anchored decision rule, score 0.694 against 0.622 in a pre-registered test
        (Sections~\ref{sec:models} and~\ref{sec:results}).
  \item \textbf{An honest record of what did not work.} More than a dozen ideas, including several that we
        expected to work, are reported with their numbers (Sections~\ref{sec:othertrials},
        \ref{sec:fakelabels} and~\ref{sec:vis}). We consider these negative results a real part of the
        contribution: they show where the difficulty of the task lies.
  \item \textbf{A reproducible evaluation.} Reproduction of all five \figsum{} baselines, a fixed split,
        out-of-fold scoring for every number, and a test set that was sealed and opened
        once (Section~\ref{sec:task}).
\end{enumerate}

\paragraph{Design constraints.} Two decisions shaped the work. First, the tool must run locally, so no
proprietary large language model is used inside the tool or to create training labels. Second, the learned
approach uses no hard distance limit: every sentence of the paper is a candidate for every figure, and distance
to the nearest citation is given to the model only as an input.
